\section{Математическое описание}

Формат bfloat16 содержит 16 бит: один бит знака $s$, восемь бит смещённой
экспоненты $E$ и семь бит дробной части мантиссы $M$. Для нормализованного числа
его значение определяется выражением

\begin{equation}
    x = (-1)^s \left(1 + M\right) 2^{E-127},
    \label{eq:normal}
\end{equation}

где

\begin{equation}
    M = \sum_{i=1}^{7} m_i 2^{-i}, \qquad m_i \in \{0,1\}.
    \label{eq:mantissa}
\end{equation}

Экспонента хранится со смещением 127, поэтому связь истинной экспоненты $e$
с её полем имеет вид

\begin{equation}
    E = e + 127.
\end{equation}

Ведущая единица нормализованного числа не сохраняется. Например,

\[
    13{,}25_{10} = 1101{,}01_2
    = 1{,}10101_2 \cdot 2^3.
\]

Следовательно, истинная экспонента равна $3$, поле экспоненты содержит
$3+127=130=10000010_2$, а поле мантиссы после дополнения нулями имеет вид
$1010100_2$.

Для преобразования десятичного числа в bfloat16 его модуль разделяется на целую
и дробную части. Целая часть переводится в двоичную систему последовательным
делением на два, дробная --- последовательным умножением на два. Полученная
двоичная запись нормализуется к виду $1{,}M\cdot 2^e$. После этого мантисса
округляется до семи бит, к экспоненте прибавляется смещение 127 и формируются
три поля результата.

При обратном преобразовании поля извлекаются из 16-битной строки. Для обычного
числа применяется формула~\eqref{eq:normal}. Значения экспоненты $0$ и $255$
зарезервированы:

\begin{itemize}
    \item $E=0$, $M=0$ --- ноль;
    \item $E=0$, $M\ne0$ --- субнормальное число
          $x=(-1)^s M 2^{1-127}$;
    \item $E=255$, $M=0$ --- положительная или отрицательная бесконечность;
    \item $E=255$, $M\ne0$ --- значение NaN.
\end{itemize}

Канонический quiet NaN в спецификации RISC-V имеет битовый шаблон
$\mathtt{0x7FC0}$, то есть знак $0$, экспоненту $\mathtt{0xFF}$ и мантиссу
$\mathtt{0x40}$.

\newpage
